\documentclass[11pt,a4paper]{article}
\usepackage[margin=1.9cm]{geometry}
\usepackage{fontspec,polyglossia}
\setdefaultlanguage{greek}\setotherlanguage{english}
\setmainfont{DejaVu Serif}\newfontfamily\greekfont{DejaVu Serif}
\setmonofont{DejaVu Sans Mono}\newfontfamily\greekfonttt{DejaVu Sans Mono}
\usepackage{amsmath,amssymb,amsthm,enumitem,comment,fancyhdr,hyperref}
\hypersetup{unicode,hidelinks}

\theoremstyle{definition}\newtheorem{question}{Άσκηση}
\newcommand{\solution}{\par\medskip\noindent\textbf{Λύση.}\par}
\newcommand{\eps}{\varepsilon}
\setlist[enumerate]{label=(\alph*),leftmargin=*}
\setlength{\parindent}{0pt}\setlength{\parskip}{6pt}
\pagestyle{fancy}\fancyhf{}\fancyfoot[L]{\small ΗΥ280 (2026) -- Τσουρακάκης}\fancyfoot[R]{\thepage}\renewcommand{\headrulewidth}{0pt}
\begin{document}
\begin{center}\rule{\linewidth}{.4pt}\\[3pt]{\large\bfseries ΗΥ 280 - Θεωρία Υπολογισμού}\\
Ασκήσεις Διαλέξεων 11, 12 - Λήμμα Άντλησης και Myhill--Nerode\\
Χαράλαμπος Ε. Τσουρακάκης\\\rule{\linewidth}{.4pt}\end{center}
\textbf{Οδηγίες.} Όλες οι ασκήσεις είναι παραδοτέες.

\begin{question}[Κλειστότητα και γινόμενο αυτομάτων]
Έστω $L_1,L_2\subseteq\{0,1\}^*$ κανονικές γλώσσες.
\begin{enumerate}
\item Αποδείξτε με κατασκευή DFA ότι η συμμετρική διαφορά $L_1\triangle L_2$ είναι κανονική.
\item Αν $L_1$ είναι οι λέξεις με άρτιο πλήθος $0$ και $L_2$ οι λέξεις με άρτιο πλήθος $1$, κατασκευάστε DFA για τη συμμετρική διαφορά (πίνακας μεταβάσεων αρκεί).
\end{enumerate}\end{question}

\begin{question}[Λήμμα άντλησης: αυστηρή ανισότητα]
Αποδείξτε ότι $A=\{0^i1^j:i,j\geq0,\ i>j\}$ δεν είναι κανονική, χρησιμοποιώντας το λήμμα άντλησης. Να αντιμετωπίσετε \emph{κάθε} νόμιμη διάσπαση.
\end{question}
 

\begin{question}[Λήμμα άντλησης: τέλειοι κύβοι]
Έστω $B=\{1^{n^3}:n\geq0\}$. Δείξτε ότι η $B$ δεν είναι κανονική. \textit{Υπόδειξη:} συγκρίνετε $p^3+t$ με $(p+1)^3$.
\end{question}

\begin{question}[Αποφασισιμότητα για DFA]
Έστω DFA $M=(Q,\Sigma,\delta,q_0,F)$ με $n=|Q|$.
\begin{enumerate}
\item Περιγράψτε αλγόριθμο προσβασιμότητας που αποφασίζει αν $L(M)=\varnothing$.
\item Περιγράψτε αλγόριθμο που αποφασίζει αν $L(M)$ είναι άπειρη, χρησιμοποιώντας κύκλους στον γράφο μεταβάσεων. Αιτιολογήστε την ορθότητα και τον τερματισμό.
\item Δώστε παράδειγμα DFA με κύκλο που αναγνωρίζει πεπερασμένη γλώσσα.
\end{enumerate}\end{question}
 

\textbf{Υπενθύμιση Myhill--Nerode.} Για $L\subseteq\Sigma^*$ ορίζουμε $x\equiv_L y$ όταν για κάθε $z\in\Sigma^*$ ισχύει $xz\in L\iff yz\in L$. Ένα επίθημα $z$ που κάνει ακριβώς μία από τις $xz,yz$ αποδεκτή \emph{διακρίνει} τις $x,y$. Η $L$ είναι κανονική αν και μόνο αν η σχέση έχει πεπερασμένο πλήθος κλάσεων. Ο αριθμός κλάσεων ισούται με το πλήθος καταστάσεων του ελάχιστου DFA.

\begin{question}$\star$
Έστω $L=\{w\in\{0,1\}^*:w\text{ τελειώνει σε }01\}$.
\begin{enumerate}
\item Δείξτε ότι οι $\eps$, $0$, $01$ ανήκουν σε διαφορετικές κλάσεις, δίνοντας διακριτικά επιθήματα.
\item Κατασκευάστε DFA με τρεις καταστάσεις για τη $L$ και αποδείξτε ότι είναι ελάχιστο.
\end{enumerate}\end{question}
 

\begin{question} $\star\star\star$
Για κάθε σταθερό $k\geq1$, έστω
\[L_k=\{w\in\{0,1\}^*:|w|\geq k\text{ και το }k\text{-οστό σύμβολο από το τέλος είναι }1\}.\]
\begin{enumerate}
\item Δείξτε ότι οι $2^k$ δυαδικές λέξεις μήκους $k$ είναι ανά δύο μη ισοδύναμες. Αν δύο λέξεις διαφέρουν στη θέση $j$ από αριστερά, δοκιμάστε επίθημα μήκους $j-1$.
\item Κατασκευάστε DFA με $2^k$ καταστάσεις για τη $L_k$. \textit{Υπόδειξη:} αποθηκεύστε τα τελευταία $k$ σύμβολα, συμπληρώνοντας αρχικά με $0$.
\item Συμπεράνετε το ακριβές πλήθος καταστάσεων του ελάχιστου DFA. Γιατί δεν αντιφάσκει αυτό με την ύπαρξη μικρού NFA για την ίδια γλώσσα;
\end{enumerate}\end{question}

\end{document}
